% Generated draft: requires mathematical and source review.
\documentclass[11pt]{article}
\usepackage[margin=25mm]{geometry}
\usepackage{fontspec}
\setmainfont{Latin Modern Roman}
\usepackage{amsmath,amssymb,mathtools,unicode-math}
\setmathfont{Latin Modern Math}
\usepackage{enumitem,needspace,xurl,hyperref}
\setlength{\parindent}{0pt}
\setlength{\parskip}{5pt}
\setlength{\emergencystretch}{3em}
\newcommand{\N}{\mathbb{N}}
\newcommand{\Z}{\mathbb{Z}}
\newcommand{\Q}{\mathbb{Q}}
\newcommand{\R}{\mathbb{R}}
\newcommand{\C}{\mathbb{C}}
\newcommand{\F}{\mathbb{F}}
\newcommand{\A}{\mathbb{A}}
\newcommand{\PP}{\mathbb{P}}
\newcommand{\E}{\mathbb{E}}
\newcommand{\eps}{\varepsilon}
\newcommand{\dd}{\,\mathrm{d}}
\begin{document}
\section*{UnsolvedMath ID: 30004180}
\textbf{Source problem code:} catalogue entry
\subsection{Definitions and mathematical statement}
Let \(D\) denote the disk appearing in the source, and let \(f\colon D\to D\) be a map in the source's class of dissipative disk diffeomorphisms. The precise meaning of this class remains unresolved: the record specifies neither the regularity, the open or closed disk convention, boundary conditions, whether the map is onto or is an embedding into the disk, nor the meaning of dissipativity.

Write \(\Z_{>0}=\{1,2,\ldots\}\) and \(\Z_{\geq0}=\{0,1,2,\ldots\}\). Define the set of least periods by
\[
\operatorname{Per}(f)=\bigl\{p\in\Z_{>0}:\exists x\in D,\ f^p(x)=x\ \text{and}\ f^j(x)\ne x\ \text{for every }1\leq j<p\bigr\},
\]
where \(f^j\) is the \(j\)-fold iterate of \(f\).

The supplied question asks whether every such \(f\) satisfies both
\[
\exists n_0\in\Z_{>0},\qquad
\operatorname{Per}(f)\subseteq
\{n2^k:n\in\Z_{>0},\ n\leq n_0,\ k\in\Z_{\geq0}\},
\]
and
\[
\operatorname{Per}(f)\text{ is infinite}
\quad\Longrightarrow\quad
\exists n\in\Z_{>0},\qquad
\{n2^k:k\in\Z_{\geq0}\}\subseteq\operatorname{Per}(f).
\]
Here \(n_0\) and \(n\) may depend on \(f\). The exponent convention is made explicit as \(k\geq0\); the source uses \(\mathbb N\) without defining it. Allowing or excluding zero does not change these existential assertions after adjusting the positive multipliers or their bound.

\textit{Unresolved scope.} No entropy hypothesis is included in either supplied formulation. The cited literature is labeled ``zero entropy,'' so the relationship between its hypotheses and the displayed question requires source review. No zero-entropy or mild-dissipativity assumption has been added to the question. Until the intended class of maps is specified, this is a conditional formalization of the supplied wording.
\subsection{Short English statement}
The question asks whether the least periods of every dissipative disk diffeomorphism come from finitely many sequences obtained by repeated doubling, and whether infinitely many distinct periods force one complete doubling sequence. The intended class of diffeomorphisms and any missing entropy hypothesis require clarification from the sources.
\subsection{Sources}
\begin{itemize}
\item Supplied local UnsolvedMath registry record, permanent ID \(30004180\), source problem code \texttt{catalogue entry}, titled \emph{Period Structure of Dissipative Disk Diffeomorphisms}, in the Oberwolfach Reports Open Problems collection. The record gives proposed year 2019 and supplies no individual proposer attribution.
\item Supplied report citation: \emph{Dynamische Systeme (2020)}, Oberwolfach Reports, pp.\ 1855--1917, DOI \texttt{10.4171/owr/2019/31}. \url{https://doi.org/10.4171/owr/2019/31}. The background describes this collection as drawn from workshop reports and their open-problem discussions. The 2020 title date and 2019 provenance are retained as supplied.
\item Supplied literature reference labeled ``Crovisier--Pujals--Tresser, zero entropy'': \url{https://arxiv.org/abs/2005.14278}.
\item Supplied literature reference labeled ``original OWR report'': \url{https://doi.org/10.4171/owr/2019/31}.
\item Literature assessment dated exactly \texttt{2026-08-22}: the supplied curation labels the problem partially solved and states that the period-doubling conclusion is proved for mildly dissipative disk diffeomorphisms, including the relevant H\'enon regime, while the exact conjecture for every dissipative disk diffeomorphism remains open. This is supplied dated curation, not an independently verified literature assessment. The precise hypotheses of these reported results and the meaning of the relevant H\'enon regime are not supplied.
\item The registry labels the statement \texttt{corrected\_verified} and describes its changes as repairs to notation, typography, and sentence structure preserving mathematical content. This is supplied editorial metadata, not independent mathematical verification of this draft or of the stated literature status.
\end{itemize}
\end{document}
